% SKELETON — structure only. Fill each \TODO; see .deepseek/13_DOCS.md §13.2.
\documentclass[11pt]{article}

\usepackage{amsmath,amssymb,amsthm}
\usepackage{algorithm,algpseudocode}
\usepackage[margin=1in]{geometry}
\usepackage{hyperref}

\newtheorem{theorem}{Theorem}
\newtheorem{proposition}{Proposition}
\newtheorem{remark}{Remark}
\newcommand{\TODO}[1]{\textbf{[TODO: #1]}}

\title{A Unified MPC-CBF Formulation}
\author{Ali-Eimaan}
\date{\today}

\begin{document}
\maketitle

\begin{abstract}
\TODO{State the single optimisation problem that specialises --- by parameter choice alone --- to
MPC-DC, MPC-CBF with fixed decay, the CDC 2021 relaxed-decay variant, and the one-step CBF-QP.
That unification is the contribution of this note and the reason the codebase has one solver class
rather than four.}
\end{abstract}

\section{The unified problem}
\TODO{Write the OCP in full:
\begin{align*}
\min_{\mathbf{u}, \boldsymbol{\omega}} \quad & \sum_{k=0}^{N-1} \left[ \|x_{k|t} - x^{\mathrm{ref}}_{k}\|_Q^2 + \|u_{k|t}\|_R^2 + p_\omega \sum_j (\omega_{k,j} - 1)^2 \right] + \|x_{N|t} - x^{\mathrm{ref}}_N\|_{Q_f}^2 \\
\text{s.t.}\quad & x_{k+1|t} = f(x_{k|t}, u_{k|t}) \\
 & x_{0|t} = x_t,\quad x_{k|t} \in \mathcal{X},\quad u_{k|t} \in \mathcal{U} \\
 & h_j(x_{k|t}) \ge 0 \\
 & h_j(x_{k+1|t}) - h_j(x_{k|t}) \ge -\omega_{k,j}\gamma h_j(x_{k|t}),\quad k < N_{\mathrm{CBF}}
\end{align*}
Then the specialisation table: $\omega \equiv 1$ gives ACC 2021; $\omega$ free with
$\omega\gamma \le 1$ gives CDC 2021; dropping the decay row gives MPC-DC; $N = 1$ with a tracking cost
gives the CBF-QP filter.}

\section{Why the decay row is imposed along the prediction}
\TODO{Contrast imposing the DCBF condition only at $k=0$ (myopic, equivalent to the QP filter) with
imposing it for $k < N_{\mathrm{CBF}}$. Explain the mechanism by which the horizon buys feasibility:
the optimiser can trade barrier decay across stages instead of paying it all immediately. Reference the
empirical curve in \texttt{analysis/feasibility\_recovery\_study.ipynb}.}

\section{Choice of $N_{\mathrm{CBF}}$}
\TODO{Constraint count is $N_{\mathrm{CBF}} \times n_{\mathrm{obs}}$; give the resulting QP size and
the measured solve-time scaling. Justify the default in \texttt{config/mpc\_cbf\_params.yaml}.}

\section{Feasibility}
\TODO{This section must be honest. State that neither ACC 2021 nor this implementation provides a
recursive-feasibility certificate; give the standard sufficient conditions (terminal control-invariant
set inside $\mathcal{C}$, or a sufficiently long horizon), state that they are \emph{not} imposed here,
and point at the empirical study. A reviewer who finds this section overclaiming discounts the whole
repository.}

\begin{proposition}[Closed-loop safety given per-step feasibility]
\TODO{If the problem is feasible at every $t$ and $x_0 \in \mathcal{C}$, then $x_t \in \mathcal{C}$ for
all $t$. Note precisely how weak this is: it assumes exactly the thing that is not guaranteed.}
\end{proposition}

\section{Numerical formulation}
\TODO{Multiple shooting; Gauss-Newton Hessian; the barrier row's Jacobian in closed form (it is affine,
so give it); why the DISCRETE integrator type is mandatory --- an ERK integrator would enforce the
condition on an interpolated state, changing its meaning.}

\begin{algorithm}
\caption{One control step}
\begin{algorithmic}[1]
\State \TODO{measure $x_t$; prune and propagate obstacles; push parameters; warm start from the shifted
previous solution; solve; extract $u_{0|t}$; shift for the next warm start; publish diagnostics.}
\end{algorithmic}
\end{algorithm}

\section{Complexity and timing}
\TODO{Report measured solve times as a function of $N$, $n_{\mathrm{obs}}$, and SQP vs RTI, from the
\texttt{SolveMeetsRealTimeBudget} gtest. Table, with the CPU named.}

\end{document}
